\section{Dar cambio en la tienda}
Considere el problema de formar una cantidad de dinero $V$ usando monedas de
denominaciones
\( \mathbf{d}= d_1, \ldots, d_n  \) de las cuales hay cantidades
\( \mathbf{q}= q_1, \ldots, q_n  \).
El programa debe imprimir
\begin{enumerate}
	\item[Caso 1:] \texttt{IMPOSSIBLE} Si no es posible formar $V$ con las
		monedas disponibles
	\item[Caso 2:] La secuencia \texttt{k1} \texttt{k2} \ldots \texttt{kn}
		de cantidades (posiblemente 0 para algunas monedas) que su programa
		encuentra para formar el valor $V$. \textbf{Ejemplo:} Si $V= 30$,
		$\mathbf{d}= 1, 5, 10, 25$ y $\mathbf{q}= 20, 0, 5, 5$ Entonces se
		imprime
		\begin{itemize}
			\item[] \texttt{0 0 3 0}.
		\end{itemize}
		La solución debe ser \textbf{óptima} en el sentido de que no hay
		soluciones que usen menos monedas.
\end{enumerate}
\[
	\begin{array}{l}
		\mbox{Ejemplo 1:}\\
		\mathbf{d}= 15, 37, 162, 170, 177, 184, 185, 192\\ 
		\mathbf{q}= 20, 20, 19,  18,  19,  11,  10,  20\\
		 \\
		\mbox{Ejemplo 2:}\\
		\mathbf{d}= 7,  68, 98, 110, 127, 140, 142, 162\\ 
		\mathbf{q}= 10, 14, 19, 17,  12,  12,  16,  18
	\end{array}
\]
Ejemplos de valores problema:
\[ V= 520, 795, 156, 269, 327, 592, 413, 850, 329, 396.\]
